\chapter*{Course Map}
\phantomsection%
\addcontentsline{toc}{chapter}{Course Map}
\markboth{Course Map}{}

The course progresses from searching a graph to learning from data. Four stages build on the toolkit in Chapter~1, and every lab follows the same five steps: define the task, validate the data, run a controlled experiment, interpret the results, and submit the executed notebook with its generated evidence.

\begin{utpfigure}
\newcommand{\coursemapcard}[3]{%
  \parbox[c][1.52cm][t]{3.35cm}{\raggedright
    {\scriptsize\bfseries\color{UTPTeal}\strut #1}\par
    \vspace{0.02em}%
    {\bfseries\strut #2}\par
    \vspace{0.08em}%
    {\scriptsize\color{UTPMuted!80!UTPNavy}\strut #3\strut}%
  }%
}
\begin{tikzpicture}[
  node distance=0.58cm and 0.66cm,
  card/.style={draw=UTPBlue,fill=UTPPaleBlue,rounded corners=2mm,
    minimum width=3.99cm,minimum height=1.88cm,align=left,
    inner xsep=3.2mm,inner ysep=1.6mm,font=\small\color{UTPNavy}},
  finalcard/.style={card,draw=UTPTeal,fill=UTPPaleTeal},
  arr/.style={-{Stealth[length=2.25mm,width=1.55mm]},line width=0.9pt,
    draw=UTPTeal,shorten >=0.8pt,shorten <=0.8pt}
]
\node[card] (a) {\coursemapcard{STAGE 01}{Route search}{BFS, DFS, UCS, and A$^\star$ on airline data}};
\node[card,right=of a] (b) {\coursemapcard{STAGE 02}{Reasoning}{Fuzzy logic, ANFIS, and robot sensing}};
\node[card,below=of b] (c) {\coursemapcard{STAGE 03}{Language}{Handwriting OCR and sentiment}};
\node[finalcard,left=of c] (d) {\coursemapcard{STAGE 04}{Optimization}{Genetic-algorithm regression}};
\draw[arr] (a) -- (b);
\draw[arr] (b) -- (c);
\draw[arr] (c) -- (d);
\end{tikzpicture}
\end{utpfigure}

\begin{utptable}
\begin{tabularx}{\linewidth}{@{}>{\bfseries\leavevmode\color{UTPNavy}\centering\arraybackslash}m{0.08\linewidth}m{0.31\linewidth}X@{}}
\toprule
\rowcolor{UTPPaleBlue!92}
\utpth{Stage} & \utpth{Focus} & \utpth{Key question} \\
\midrule
1 & Uninformed and informed search & How is a good route found, and what does finding it cost? \\
2 & Fuzzy systems and robot sensing & How can imprecise readings drive a sound decision? \\
3 & OCR and text classification & How is text turned into evidence a model can use? \\
4 & Evolutionary optimization & How can a population of candidates fit a model? \\
\bottomrule
\end{tabularx}
\end{utptable}

\section*{Datasets and what each stage builds}
Every stage works on one or two public Kaggle datasets, listed with download links in Chapter~1.

\begin{utptable}
\begin{tabularx}{\linewidth}{@{}>{\bfseries\leavevmode\color{UTPNavy}\centering\arraybackslash}m{0.08\linewidth}>{\centering\arraybackslash}m{0.09\linewidth}m{0.30\linewidth}X@{}}
\toprule
\rowcolor{UTPPaleBlue!92}
\utpth{Stage} & \utpth{Labs} & \utpth{Dataset} & \utpth{What you build} \\
\midrule
1 & 1--3 & OpenFlights routes and airports & BFS, DFS, UCS, GBFS, and A$^\star$, then a tested command-line toolkit. \\
2 & 4--5 & UCI Air Quality; SCITOS wall-following sensor readings & Fuzzy and ANFIS models, and an interpretable sensor-to-action controller. \\
3 & 6--7 & Doctor Handwriting Recognition; Twitter Airline Sentiment & A handwriting OCR evaluation and a three-class airline sentiment benchmark. \\
4 & 8 & Medical Cost Personal Dataset & A repeated-run GA regression compared with least squares. \\
\bottomrule
\end{tabularx}
\end{utptable}
